\chapter{Conclusiones y Trabajos Futuros}

\section{Conclusiones}

El presente proyecto ha abordado el desarrollo de un sistema de detección de lenguaje ofensivo en español, utilizando modelos de lenguaje basados en transformers, con el fin de responder a una necesidad creciente en los entornos digitales actuales. A través del uso del corpus OffendES y del modelo BETO, se logró construir una arquitectura de clasificación que permite identificar con buena precisión comentarios ofensivos dentro de contextos reales de comunicación, como los que se producen en redes sociales o medios digitales.

Durante el desarrollo se evidenció que, si bien existen herramientas avanzadas utilizadas por grandes plataformas internacionales, muchas de estas soluciones no están disponibles para terceros o no contemplan adecuadamente el idioma español. Esta situación reafirma la relevancia y oportunidad del presente trabajo, especialmente al enfocarse en un mercado objetivo como el de los medios de comunicación digitales, los cuales requieren soluciones ágiles y confiables para gestionar comentarios del público en sitios web y transmisiones en vivo.

El sistema propuesto no solo permite clasificar comentarios de manera eficiente, sino que también es adaptable a diferentes plataformas y escalable según la demanda. Se ha demostrado la viabilidad técnica del modelo y se ha propuesto un esquema de despliegue funcional, además de una evaluación económica que sustenta el potencial de este proyecto como servicio en el mercado.

\subsection{Trabajos Futuros}

A partir de los resultados obtenidos, se plantean las siguientes líneas de trabajo para fortalecer y extender el alcance del sistema:

\begin{itemize}
    \item \textbf{Ampliación del corpus:} Incorporar nuevos datasets multitemáticos o propios, que reflejen mejor la variabilidad del lenguaje ofensivo en otras plataformas como foros, chats en vivo, videojuegos y transmisiones.

    \item \textbf{Clasificación multicategoría:} Extender la arquitectura del modelo para no solo detectar la presencia de lenguaje ofensivo, sino también categorizarlo según su tipo (insultos, discurso de odio, violencia de género, etc.).
\end{itemize}

% =========================
% Página 68
% =========================

\vspace{0.5cm}

\begin{itemize}
    \item \textbf{Implementación de una interfaz web o API pública:} Permitir el uso del modelo por parte de terceros a través de una API REST o una extensión de navegador, facilitando la integración con otros servicios.

    \item \textbf{Adaptación a tiempo real:} Optimizar el rendimiento del modelo para implementaciones que requieran respuesta inmediata, como transmisiones en vivo o bots de moderación en plataformas como Discord, Telegram o transmisiones televisivas.

    \item \textbf{Evaluación en escenarios reales:} Realizar pruebas con medios de comunicación locales u organizaciones interesadas en moderar sus espacios de participación, para validar el sistema en condiciones reales de operación.

    \item \textbf{Refinamiento lingüístico y cultural:} Mejorar la detección tomando en cuenta variaciones regionales del español y contextos culturales que puedan afectar la interpretación de lo que es ofensivo.
\end{itemize}

% =========================
% Página 69
% =========================

